% ==============================================================
% Hilbert Spaces and Qubits
% ==============================================================

\section{Hilbert Spaces and Qubits}

\subsection{Qubits and Dirac notation}

% TODO

\subsection{Tensor products and multi-qubit systems}

\begin{definition}
The \textbf{tensor product} (written $\otimes$) combines two quantum systems into a joint system. If qubit 1 is in state $|a\rangle$ and qubit 2 is in state $|b\rangle$, their combined state is:
\[
|a\rangle \otimes |b\rangle
\]
This is usually abbreviated by juxtaposition: $|a\rangle|b\rangle = |ab\rangle$.
\end{definition}

\begin{example}[Two-qubit states]
Using the standard basis vectors $|0\rangle = \begin{bmatrix}1\\0\end{bmatrix}$ and $|1\rangle = \begin{bmatrix}0\\1\end{bmatrix}$, the tensor product is computed as the Kronecker product — multiply the scalar entries of the first vector by the entire second vector:
\[
|0\rangle \otimes |1\rangle = \begin{bmatrix}1\\0\end{bmatrix} \otimes \begin{bmatrix}0\\1\end{bmatrix} = \begin{bmatrix}1\cdot\begin{bmatrix}0\\1\end{bmatrix}\\[4pt] 0\cdot\begin{bmatrix}0\\1\end{bmatrix}\end{bmatrix} = \begin{bmatrix}0\\1\\0\\0\end{bmatrix}
\]
This 4-dimensional vector is the 2-qubit basis state $|01\rangle$. Similarly, $|1\rangle \otimes |0\rangle = |10\rangle = \begin{bmatrix}0\\0\\1\\0\end{bmatrix}$.

An $n$-qubit system lives in a $2^n$-dimensional Hilbert space — the state space grows exponentially with the number of qubits.
\end{example}

\begin{remark}[Tensor product for gates]
The $\otimes$ notation also applies to quantum gates. Writing $H^{\otimes 2} = H \otimes H$ means ``apply the Hadamard gate independently to each of two qubits''. More generally, $U^{\otimes n}$ applies gate $U$ in parallel to all $n$ qubits.
\end{remark}

\subsection{Inner products and measurement}

% TODO

\subsection{Circled operator notation}

\begin{remark}[Reference: circled operators]
Several ``circled'' binary operators appear in mathematics and quantum computing. Their meaning is context-dependent:

\begin{center}
\begin{tabular}{c|l|l}
Symbol & Name & Common meaning \\
\hline
$\oplus$ & Circled plus & XOR / addition mod 2 / direct sum \\
$\otimes$ & Circled times & Tensor product / Kronecker product \\
$\ominus$ & Circled minus & Formal difference / subtraction-like operation \\
$\oslash$ & Circled slash & Quotient-like / author-defined division operation \\
$\odot$ & Circled dot & Elementwise (Hadamard) product / composition \\
$\circ$ & Small circle & Function composition: $(f \circ g)(x) = f(g(x))$ \\
$\bigoplus$ & Big circled plus & Direct sum over many terms: $\bigoplus_{i=1}^n V_i$ \\
$\bigotimes$ & Big circled times & Tensor product over many terms: $\bigotimes_{i=1}^n |q_i\rangle$ \\
\end{tabular}
\end{center}

\textbf{Quick intuition:}
\begin{itemize}
    \item $\oplus$ — special addition (bit-flipping arithmetic)
    \item $\otimes$ — gluing systems together (combining spaces)
    \item $\ominus$ — special subtraction
    \item $\oslash$ — special division
    \item $\odot$ — special elementwise multiplication
    \item $\circ$ — function composition (not a circled symbol, but related in appearance)
\end{itemize}
Unlike the standard $+, -, \times$, circled symbols are \emph{not universal} — always check the author's definition when encountering them in a new context.
\end{remark}
